\chapter{Conclusions}
\label{ch:implications}

This chapter interprets the empirical findings from Chapter~\ref{ch:results}, answers the three research questions, and derives practical implications. The discussion synthesises quantitative evidence from 230 applications generated by 10 large language models and analysed by 22 automated tools across four containerised services.

%==============================================================================
\section{Addressing the research questions}
\label{sec:research_questions}
%==============================================================================

\subsection{RQ1: Security posture across models}

\textit{Which vulnerability classes (CWE/OWASP) most frequently appear in LLM-generated full-stack web applications, and how do static and dynamic security results vary across the evaluated models?}

The security landscape of LLM-generated code is dominated by configuration-level weaknesses rather than logic-level vulnerabilities. Dynamic analysis via OWASP ZAP identified 3{,}461 alerts across 295 successful scans, with the five most frequent CWEs being:

\begin{enumerate}
    \item \textbf{CWE-693} (Protection Mechanism Failure): 1{,}132 alerts---primarily missing Content Security Policy headers, indicating that no model consistently generates CSP configurations without explicit prompting.
    \item \textbf{CWE-497} (Exposure of Sensitive System Information): 888 alerts---server version disclosure via HTTP response headers, a Flask default that models do not override.
    \item \textbf{CWE-264} (Permissions, Privileges, and Access Controls): 528 alerts---cross-domain misconfiguration stemming from the CORS setup in the scaffolding and model-generated route handling.
    \item \textbf{CWE-1104} (Use of Unmaintained Third-Party Components): 295 alerts---reflecting dependency age rather than model behaviour.
    \item \textbf{CWE-1021} (Improper Restriction of Rendered UI Layers): 292 alerts---missing anti-clickjacking headers (X-Frame-Options), uniformly absent across all models.
\end{enumerate}

Critically, ZAP identified zero high-severity alerts across all 230 applications. All dynamic findings were medium or low severity, indicating that LLM-generated applications do not exhibit the critical runtime vulnerabilities (SQL injection, remote code execution, authentication bypass) that might be expected from automated code generation.

Static analysis tells a complementary story. Bandit identified 924 security-relevant findings (222 high, 702 low), averaging 4.0 per application. The high-severity Bandit findings were concentrated in GLM-4.7 (6.7/app) and DeepSeek R1 (4.5/app), while GPT-4o Mini produced the fewest (1.4/app). Semgrep detected 1{,}690 findings, all classified as low severity, confirming that LLM-generated code contains insecure patterns (e.g., missing input validation, broad exception handling) but not exploitable vulnerabilities in isolation.

Inter-model variation in security findings is substantial. ZAP alerts per successful scan ranged from 11.8 (GPT-5.2 Codex and Gemini 3 Flash, both at 40 scans) to 11.8 (Claude 4.5 Sonnet at 93 scans), showing remarkable uniformity when normalised by scan count---the apparent variation in raw totals (96--1{,}101) is almost entirely explained by differing deployment success rates rather than inherent security differences. Models with 0\% deployment (Mistral Small 3.1, Llama 3.1 405B) produced zero ZAP alerts simply because no containers were available to scan, representing a measurement gap rather than superior security.

\textbf{Answer to RQ1:} The most frequent vulnerability classes are configuration-level HTTP security header omissions (CWE-693, CWE-497, CWE-1021) rather than logic-level code vulnerabilities. No high-severity dynamic vulnerabilities were detected. Static security tools (Bandit, Semgrep) find low-to-moderate numbers of security-relevant patterns, with inter-model variation of up to 4.8$\times$ in Bandit findings per application, driven primarily by code volume differences rather than fundamentally different security practices.

\subsection{RQ2: Correctness and reliability vs.\ security}

\textit{How do functional compliance (requirements coverage and API conformance) and build/run success relate to observed security findings, and are there systematic trade-offs between correctness and security?}

The study reveals a nuanced relationship between correctness and security rather than a simple trade-off. Requirements compliance ranged from 69.9\% (Gemini 3 Pro) to 91.0\% (Mistral Small 3.1), with a cross-model mean of 81.6\%. Deployment success---the proportion of applications that built, started, and responded to HTTP probes---showed even wider variation, from 0\% (Mistral Small 3.1, Llama 3.1 405B) to 100\% (GPT-5.2 Codex, Gemini 3 Flash).

The most striking finding is the \textit{compliance-deployment paradox}: Mistral Small 3.1 achieves the highest requirements compliance (91.0\%) yet fails to produce any deployable containers. Its generated code correctly implements specified features at the source level but contains structural or dependency issues that prevent successful containerised deployment. Conversely, GPT-5.2 Codex achieves 100\% deployment with moderate compliance (83.2\%), suggesting that deployment success depends more on infrastructure integration competence (correct imports, dependency declarations, configuration files) than on feature completeness.

Defect density (D/kLOC) does not show a simple inverse relationship with compliance. GPT-5.2 Codex exhibits both the lowest defect density (76.8 D/kLOC) and the highest code volume (2{,}707 LOC/App), indicating that its additional output represents well-structured code rather than inflated boilerplate. Claude 4.5 Sonnet, producing comparable volume (2{,}645 LOC/App), has nearly double the defect density (139.3 D/kLOC), driven by a disproportionate number of low-severity linting findings (61\% Low, 22\% Medium, 17\% High). This suggests that Claude's verbose output style triggers more linter warnings without necessarily indicating worse security.

Security findings (Bandit + Semgrep combined) correlate weakly with compliance ($r \approx 0.2$), confirming that implementing more features does not systematically introduce more security vulnerabilities. The absence of high-severity ZAP findings across all models, regardless of compliance level, further supports this conclusion.

\textbf{Answer to RQ2:} There is no systematic trade-off between correctness and security. High compliance does not predict high vulnerability counts, and deployment success is orthogonal to requirements coverage. The critical distinction is between \textit{feature-level correctness} (compliance) and \textit{infrastructure-level correctness} (deployment), which appear to be independent model capabilities. GPT-5.2 Codex uniquely excels at both, while Mistral Small 3.1 demonstrates that source-level feature completeness does not guarantee operational readiness.

\subsection{RQ3: Code quality, maintainability, and performance}

\textit{How do engineering-quality indicators (linting, type coverage, complexity, dead code) and performance under load differ across models, and how do these properties correlate with security outcomes?}

Code quality metrics reveal significant inter-model variation. Pylint findings per application ranged from 21.0 (Llama 3.1 405B) to 107.7 (Claude 4.5 Sonnet), a 5.1$\times$ difference. Ruff findings showed similar spread: 9.7 (Llama 3.1 405B) to 128.8 (Claude 4.5 Sonnet). However, these absolute differences are strongly confounded by code volume---models producing more code naturally trigger more linting findings. When normalised as I/100LOC (Issues per 100 Lines of Code), the spread narrows considerably: 7.7 (GPT-5.2 Codex) to 13.9 (Claude 4.5 Sonnet), a 1.8$\times$ ratio.

The correlation analysis from Section~\ref{sec:model_rankings} reveals that the strongest predictor of code quality is \textit{maximum output token capacity} ($\rho = -0.77$ with D/kLOC). Models with higher output limits produce code with lower defect density, likely because larger output windows allow more complete, self-consistent implementations without truncation artefacts. Context window size shows weaker correlation with quality ($\rho = -0.26$ with D/kLOC), while pricing shows negligible correlation ($\rho = -0.09$), suggesting that more expensive models do not automatically produce higher-quality code.

Performance under load was remarkably uniform across successfully deployed models. Backend throughput ranged from 330.6 RPS (Gemini 3 Flash) to 388.5 RPS (Qwen3 Coder+), a mere 17.5\% spread. Median response times clustered tightly between 2.57\,ms and 2.66\,ms for most models, with outliers only in Gemini 3 Flash (2.66\,ms median but 14.4\,ms maximum) and GLM-4.7 (2.65\,ms median but elevated mean of 4.12\,ms due to occasional slow requests). This uniformity is expected: all applications share the same Flask scaffolding, SQLite backend, and Docker resource limits, so performance differences reflect the efficiency of model-generated route handlers rather than architectural choices.

The relationship between code quality and security is weak. Pylint and Ruff findings (code style, conventions) show minimal correlation with Bandit findings (security patterns), confirming that code quality tools and security tools measure largely orthogonal dimensions. A well-linted application is not necessarily more secure, and a high-Bandit application may still pass quality checks. This supports the argument for multi-modal analysis: no single tool category captures the full quality-security spectrum.

\textbf{Answer to RQ3:} Code quality varies by up to 1.8$\times$ across models when normalised for code volume, with maximum output token capacity being the strongest predictor ($\rho = -0.77$). Performance is largely uniform (17.5\% throughput spread), determined by the shared infrastructure rather than model-generated code. Code quality and security metrics are weakly correlated, confirming that they represent independent dimensions requiring separate tool categories for comprehensive evaluation.

%==============================================================================
\section{Maturity of AI code generation}
%==============================================================================

\subsection{Readiness for production}

Assessing production readiness requires defining threshold criteria across multiple dimensions. Based on the study's metrics, a minimal readiness profile for a generated web application comprises: (1) successful deployment as a Docker container responding to HTTP probes, (2) requirements compliance above 80\%, (3) zero high-severity Bandit findings, (4) defect density below 100 D/kLOC, and (5) backend throughput above 300 RPS.

Applying these thresholds reveals that only \textbf{two models}---GPT-5.2 Codex and Gemini 3 Flash---consistently produce applications meeting all five criteria. GPT-5.2 Codex achieves 100\% deployment, 83.2\% compliance, 76.8 D/kLOC, and 383.6 RPS. Gemini 3 Flash achieves 100\% deployment, 87.7\% compliance, 107.8 D/kLOC (marginally above threshold), and 330.6 RPS. The remaining eight models fail at least one criterion, most commonly deployment success.

This finding carries practical implications: even with structured scaffolding (immutable Docker configuration, pre-built authentication, standardised dependencies), \textbf{80\% of evaluated models cannot reliably produce deployable applications}. Without scaffolding assistance, the failure rate would likely be higher. AI-generated code in its current state requires mandatory human review and automated testing before production deployment, particularly for infrastructure integration (dependency management, configuration, containerisation) where models most frequently fail.

\subsection{Complexity limits}

The study's 20 requirement templates each specified $\leq$5 backend, frontend, and admin requirements---deliberately constrained complexity to enable fair cross-model comparison. Even at this moderate complexity level, several limitations emerged:

\begin{itemize}
    \item \textbf{Deployment failures:} Two models (Mistral Small 3.1, Llama 3.1 405B) achieved 0\% deployment across all 20 templates, indicating that single-file merged generation with moderate requirements exceeds their reliable capability.
    \item \textbf{Code volume ceiling:} Llama 3.1 405B produced the least code (966 LOC/App), constrained by its 10k context window (versus 128k--1{,}048k for other models), suggesting that context capacity directly limits implementation completeness.
    \item \textbf{Consistency degradation:} Gemini 3 Pro exhibited the widest compliance variability ($\sigma = 38.1$), with individual applications ranging from 0\% to 100\% compliance, indicating that template complexity affects output stability even when mean performance is moderate.
\end{itemize}

These findings suggest that current LLMs can reliably generate simple CRUD applications with standardised scaffolding assistance but struggle with integration complexity beyond a few interacting components. Scaling to microservice architectures, real-time features, or complex state management would likely amplify the observed failure modes.

%==============================================================================
\section{Security posture of generated code}
\label{sec:security_posture}
%==============================================================================

\subsection{Default insecurity pattern}

The most pervasive security finding is what can be termed the \textit{default insecurity pattern}: LLM-generated applications inherit the insecure defaults of their frameworks without implementing security hardening. Specifically:

\begin{itemize}
    \item \textbf{Missing HTTP security headers:} No model consistently generates Content Security Policy (1{,}132 CWE-693 alerts), X-Content-Type-Options (292 alerts), or X-Frame-Options (292 alerts) headers. These are one-line Flask middleware additions that models omit because their training data predominantly contains applications without explicit security headers.
    \item \textbf{Server information disclosure:} Flask's default \texttt{Server} header exposes version information (888 CWE-497 alerts). Models do not add the \texttt{app.config['SERVER\_NAME']} suppression or middleware to strip this header.
    \item \textbf{Broad CORS configuration:} The scaffolding includes Flask-CORS with permissive defaults. Models rarely restrict CORS origins, methods, or headers beyond the scaffolding's baseline configuration.
\end{itemize}

This pattern is consistent with training data bias: LLMs learn from repositories where security configuration is often absent from example code and tutorials. The implication is that \textit{security-by-default cannot be assumed for LLM-generated code}---organisations must either encode security requirements in scaffolding (as this study partially does with authentication) or enforce post-generation security gates.

Notably, the \textit{absence} of certain vulnerability classes is also informative. Detect-secrets found zero hardcoded credentials across all 230 applications, indicating that models have learned to avoid embedding secrets in source code. Similarly, Radon identified minimal cyclomatic complexity issues (9 total findings), suggesting that LLM-generated functions are generally well-decomposed. These positive patterns likely reflect the prevalence of clean-code practices in training data.

\subsection{Dependency supply chain risk}

Dependency scanning revealed a uniform vulnerability surface across all models. Safety identified exactly 9.0 findings per application for every model (180 total per 20-app cohort), confirming that these vulnerabilities originate from the shared scaffolding's pinned dependency versions (Flask 3.0.0, SQLAlchemy 2.0.25, etc.) rather than model-generated code. Similarly, npm-audit found 1.8--2.0 findings per application uniformly, reflecting React 18.2 and Vite 5.0 dependency chains.

Pip-audit showed slightly more variation (4.5--9.5 per application across models), suggesting that some models add Python dependencies with known vulnerabilities. However, the variation is modest and driven by a small number of packages.

The dependency risk is significant in absolute terms---each application inherits approximately 11 known dependency vulnerabilities (9 Python + 2 JavaScript) before any model-generated code is considered. This \textit{inherited vulnerability baseline} must be addressed through dependency management policies (regular updates, version pinning with automated scanning) rather than model selection, since the attack surface is determined by the technology stack, not the LLM.

%==============================================================================
\section{Utility of the AutoLLM Code Analyzer platform}
%==============================================================================

\subsection{Multi-modal analysis necessity}

The study's four-service analysis architecture (static, dynamic, performance, AI) demonstrates that no single analysis modality provides comprehensive coverage. Static analysis processed all 230 applications and identified 49{,}657 findings, but these are exclusively source-code-level issues (linting violations, type errors, insecure patterns) that cannot detect runtime behaviour. Dynamic analysis (ZAP) identified 3{,}461 alerts that static tools cannot detect---HTTP header misconfigurations, cross-domain issues, and server information disclosure are only observable in running applications.

The coverage gap is quantified by service completion rates: static and AI analysis achieved 100\% completion, while dynamic and performance analysis succeeded on only 63.0\% and 63.5\% of applications respectively. This means that \textbf{37\% of applications were evaluated by only two of four analysis modalities}, creating blind spots in the security assessment. For non-deployable applications, the platform falls back to source-code-only analysis, which misses the runtime vulnerability classes that dominate the CWE findings.

Performance testing adds a third dimension that neither static nor dynamic analysis captures. The narrow throughput range (330--389 RPS) across deployed models suggests that performance is not a differentiating factor for simple CRUD applications, but the framework would be essential for more complex workloads where model-generated inefficiencies (e.g., N+1 query patterns, synchronous I/O blocking) could cause measurable degradation.

AI-based analysis (requirements compliance, code quality scoring) provides the only assessment of \textit{functional correctness}---whether the generated application actually implements the specified requirements. This dimension is orthogonal to security and quality metrics: an application can score highly on linting and pass all ZAP scans while implementing none of the required features. The 69.9\%--91.0\% compliance range demonstrates that functional correctness varies substantially and independently of other quality indicators.

The platform's containerised architecture ensures reproducible analysis: identical tool versions, configurations, and execution environments across all 230 applications. This eliminates environmental confounds that plague manual security assessments and enables fair cross-model comparison.

%==============================================================================
\section{Limitations and threats to validity}
\label{sec:limitations}

Reference is made to the methodological constraints outlined in the study design. This study focuses on Flask-based Python web applications with React frontends deployed in standardized Docker containers, limiting generalizability to other frameworks and environments. The results represent a snapshot of LLM capabilities as of June 15, 2025 and may change as models evolve. Due to time and financial constraints, the analysis relies primarily on automated tools (e.g., Bandit, OWASP ZAP) rather than comprehensive manual review. While standardized deployment improves experimental control, it may not fully reflect production environments with diverse infrastructure.

\subsection{Operational constraints and tooling variability}

Several practical factors cannot be fully standardized across providers. Token limits differ by model and gateway (e.g., OpenAI GPT-4o Mini: 16,384; Anthropic Claude 4.5 Sonnet: 64,000; Google Gemini 3: 32,768; default: 8,192), which can affect completeness when generating multi-file applications with non-trivial dependency graphs. Even when prompts are held constant, lower context ceilings can force truncation or omission of auxiliary files (e.g., configuration, tests, or security headers), creating confounds that reflect provider constraints rather than pure model capability.

In addition, a substantial portion of high-performing frontier models are closed-source and do not disclose sufficient architectural details to support fine-grained, architecture-level analysis (e.g., parameter count, MoE routing, or training corpus characteristics). This limits the interpretability of cross-model differences and reduces the ability to attribute observed behaviors to specific design choices.

Finally, parts of the analysis pipeline rely on automated code analysis, including AI-assisted components. While prior work reports that AI-enhanced code review can be effective in identifying issues and proposing fixes \cite{almeida_aicodereview_2024}, in practice such tools can be problematic: outputs may be sensitive to context and formatting, may over-report low-signal findings, or may miss issues that require deeper program understanding. Industry guidance also frames AI code analysis as a promising but still evolving capability requiring careful validation and human oversight \cite{gitlab_ai_code_analysis_2026}.

\subsection{Generation mode and output stability}

Only \texttt{GUARDED} mode is used due to \texttt{UNGUARDED} (raw) generation exhibiting a high failure rate in preliminary trials. Typical failure modes included incomplete repositories, inconsistent file structures, dependency mismatches (e.g., missing or incompatible packages), and reduced consistency across repeated generations under the same prompt. As a result, the reported outcomes primarily characterize structured generation scenarios where the pipeline enforces scaffolding and validation, and may understate the variance and brittleness observed in raw model outputs.

\subsection{Model and evaluation timing}
\label{sec:validity_timing}

Results represent a December 2025 snapshot of rapidly evolving model ecosystems. Model weights, system prompts, safety layers, default tool-use behaviors, and pricing can change on short timescales, meaning that relative rankings and observed security/quality characteristics may shift quickly. Consequently, the study should be interpreted as an empirical baseline for the specified model versions and tooling at the time of collection, with an expectation that relevance may decay as newer model releases supersede the evaluated systems.

%==============================================================================

%==============================================================================
\section{Recommendations for practice}
%==============================================================================

Based on the empirical findings, the following recommendations are directed at organisations integrating LLM code generation into their development workflows:

\begin{enumerate}
    \item \textbf{Mandate security scaffolding.} The default insecurity pattern (Section~\ref{sec:security_posture}) demonstrates that LLMs do not generate HTTP security headers, CORS restrictions, or server hardening configurations. Organisations should provide pre-configured security middleware (CSP headers, X-Frame-Options, server header suppression) as immutable scaffolding rather than expecting models to generate these protections.

    \item \textbf{Implement automated security gates.} Every generated application should pass a minimum analysis pipeline before deployment: static security scanning (Bandit + Semgrep), dependency auditing (Safety + pip-audit + npm-audit), and---for deployable applications---dynamic scanning (OWASP ZAP). The study's 100\% static analysis completion rate demonstrates that this is operationally feasible.

    \item \textbf{Select models with deployment awareness.} Model selection should prioritise deployment success rate alongside code quality metrics. The compliance-deployment paradox shows that feature-level correctness (compliance) does not predict operational readiness. GPT-5.2 Codex and Gemini 3 Flash demonstrated 100\% deployment success; organisations should benchmark candidate models on deployment reliability, not just code quality benchmarks.

    \item \textbf{Prefer models with large output windows.} The strong correlation between maximum output token capacity and defect density ($\rho = -0.77$) suggests that models with higher output limits produce more complete, lower-defect code. When cost permits, organisations should select models with $\geq$64k output token capacity.

    \item \textbf{Maintain dependency hygiene independently.} Dependency vulnerabilities (Safety: 9.0/app, npm-audit: $\sim$2.0/app) are uniform across models and originate from the technology stack, not the LLM. Organisations must implement automated dependency update pipelines (Dependabot, Renovate) regardless of which model generates the application code.

    \item \textbf{Apply multi-modal analysis.} No single tool category provides comprehensive coverage. Static analysis misses runtime configuration issues (37\% of applications lacked dynamic analysis data), while dynamic analysis cannot assess source-code quality or requirements compliance. A minimum viable pipeline should include at least static security, dependency audit, and AI-based compliance checking.

    \item \textbf{Do not assume production readiness.} Only 2 of 10 evaluated models (20\%) consistently meet minimal production-readiness thresholds across deployment, compliance, defect density, and performance. Generated code should be treated as a \textit{first draft} requiring human review, automated testing, and security hardening before production deployment.
\end{enumerate}

%==============================================================================
\section{Future research directions}
%==============================================================================

\subsection{Advanced automated remediation}

The study identifies recurring vulnerability classes (CWE-693, CWE-497, CWE-1021) that are amenable to automated remediation. Future work should investigate \textit{feedback-loop architectures} where security tool findings (in SARIF format) are fed back to the generating LLM for targeted remediation. The platform's existing SARIF integration provides a natural interface for this: Bandit and Semgrep output could be parsed into remediation prompts, creating an iterative generate-analyse-remediate cycle. The key research question is whether LLMs can reliably fix the security issues they introduce, or whether automated fixes introduce new vulnerabilities.

\subsection{Broader model and framework evaluation}

This study evaluates 10 models generating Flask/React applications. The rapid emergence of new models (e.g., GPT-5, Claude 5, Gemini 4) and specialised coding models (e.g., Cursor-based models, fine-tuned code generators) necessitates continuous evaluation. Equally important is extending the technology stack beyond Flask/React to include Django, FastAPI, Express.js, Spring Boot, and mobile frameworks, as security patterns may differ substantially across frameworks. The platform's modular architecture (interchangeable analyzer services) supports this extension with minimal modification.

\subsection{Template expansion and complexity scaling}

The current 20 templates with $\leq$5 requirements each represent moderate complexity. Future studies should systematically increase complexity along multiple axes: microservice architectures (inter-service communication, API gateways), real-time features (WebSocket, Server-Sent Events), complex state management (Redux, Zustand), ML model integration, and multi-tenant authentication. The hypothesis is that deployment success and compliance will degrade non-linearly with complexity, and identifying the inflection points would provide practical guidance on the boundary of reliable AI code generation.

\subsection{Prompt engineering for security}

The study uses standardised prompts without explicit security instructions. A controlled experiment varying prompt security directives (e.g., ``implement OWASP security headers'', ``follow secure coding guidelines'', ``add input validation to all endpoints'') would quantify the impact of security-aware prompting on vulnerability counts. Preliminary evidence from the scaffolding approach (pre-built authentication reduces auth-related vulnerabilities to zero) suggests that explicit security guidance in prompts could substantially reduce the default insecurity pattern, but this requires systematic validation across models and vulnerability classes.
